\scanpage{442}
\mspage{354}
vraie, i.e que pour toute \deleted{partie} sous-catégorie $\underline S$ de $M$ accessible dans $M$, \deleted{\unsure{$f$}} son image essentielle dans $N$ est accessible dans $N$. Dans cette question, $\mathrm{OPS}$ que $\underline S=M$, et la question revient à celle-ci : tt foncteur accessible \inserted{$f:M\unsure{\text{---}|\longrightarrow}N$} entre catégories accessibles se factorise-t-il en un foncteur (accessible) \uline{ess. surjectif}, et un foncteur accessible \uline{pl. fidèle} \diagram{\begin{tikzcd}[ampersand replacement=\&,column sep=small] M\arrow[r]\&\underline T\arrow[r]\&N \end{tikzcd}} (où $\underline T$ est elle-$=$ accessible) ?
% slanted
\margin{Il faut se borner tout au moins à ess. \unsure{$G$-surjectif|$\mathcal G$-surjectif}, mais $=$ cel\unsure{à|a} est très optimiste.}
Si oui, cela me paraîtrait peut-être le résultat le plus profond (peut-être le seul \unsure{véritablemt|véritablement} profond) en théorie des catégories.

Voici une question qui me paraît très voisine. Soient $M$ une catégorie accessible, $A$ une sous-catégorie pleine petite, Supposons que $M$ satisfasse à $L_{\pi_0}$ (p. ex. $\pi_0$ est utile pour $M$), et soit $\underline S$ la plus petite sous-catégorie str. pleine de $M$ contenant $A$, et stable par limites grdes devant $\pi_0$. \deleted{Alors} Je conjecture que $\underline S$ est accessible dans $M$.
% archive sheet number: 441
